%! Tex main = MarquezSalazarBrandon-PropuestaDeTesis.tex
\subsection{Aportación a problemas prioritarios}

Las características del modelo físico tratado en este proyecto
(alta dimensionalidad, múltiples óptimos locales, no
linealidad y una estructura matemática compleja) son las
mismas que aparecen en problemas de optimización asociados a
diversos Programas Nacionales Estratégicos (PRONACES) de
CONAHCYT. Por ejemplo, el diseño de redes eléctricas
inteligentes y la integración de energías renovables
(PRONACE 11, Transición energética y cambio climático), así
como el manejo sustentable de recursos naturales y la
conservación de ecosistemas (PRONACE 3, Sistemas
socioambientales y sustentabilidad), requieren explorar
espacios de soluciones extensos, escapar de óptimos locales y
encontrar configuraciones estables. La implementación y
comparación de distintas familias de métodos inspirados en la
naturaleza (NIM) sobre un sistema de alta complejidad
contribuye a consolidar el conocimiento metodológico necesario
para abordar esta clase de problemas. Las herramientas
computacionales y el análisis comparativo desarrollados en
este proyecto aportan así elementos transferibles a dichos
contextos, no de manera directa, sino como parte del acervo
metodológico disponible para la investigación aplicada en
áreas prioritarias.
